\section*{Bab \ref{ch:g12:trigo} --- Fungsi Trigonometri}

\begin{solution}{exo:g12:trigo:1}
$\cos\frac{2\pi}{3} = \cos\left(\pi - \frac\pi3\right) = -\cos\frac\pi3
= -\frac12$.

$\sin\frac{5\pi}{6} = \sin\left(\pi - \frac\pi6\right) = \sin\frac\pi6
= \frac12$.

$\cos\frac{7\pi}{4} = \cos\left(2\pi - \frac{\pi}{4}\right)
= \cos\left(-\frac\pi4\right) = \frac{\sqrt2}{2}$.

$\sin\left(-\frac\pi3\right) = -\sin\frac\pi3 = -\frac{\sqrt3}{2}$.
\end{solution}

\begin{solution}{exo:g12:trigo:2}
$f'(x) = 2\sin x\cos x = \sin 2x$ (aturan rantai).

$g'(x) = -3\sin(3x) + \sin x + x\cos x$ (aturan rantai dan aturan hasil kali).

Aturan hasil bagi:
\[
h'(x) = \frac{\cos x\,(2 + \cos x) - \sin x\,(-\sin x)}{(2+\cos x)^2}
= \frac{2\cos x + \cos^2 x + \sin^2 x}{(2+\cos x)^2}
= \frac{2\cos x + 1}{(2+\cos x)^2}.
\]
\end{solution}

\begin{solution}{exo:g12:trigo:3}
\emph{(a)} Satu penyelesaian khusus $\sin x = \frac{\sqrt3}{2}$ adalah
$\frac\pi3$. Penyelesaian umumnya: $x = \frac\pi3 + 2k\pi$ atau
$x = \pi - \frac\pi3 + 2k\pi = \frac{2\pi}{3} + 2k\pi$. Di
$\intoc{-\pi}{\pi}$: $x \in \left\{\frac\pi3, \frac{2\pi}{3}\right\}$.

\emph{(b)} $\cos 2x = \frac{\sqrt2}{2} = \cos\frac\pi4$ memberi
$2x = \pm\frac\pi4 + 2k\pi$, jadi $x = \pm\frac\pi8 + k\pi$. Di
$\intoc{-\pi}{\pi}$:
$x \in \left\{-\frac{7\pi}{8}, -\frac\pi8, \frac\pi8, \frac{7\pi}{8}\right\}$.
\end{solution}

\begin{solution}{exo:g12:trigo:4}
Dengan $a = \frac\pi3$, $b = \frac\pi4$:
\[
\cos\frac{\pi}{12} = \cos a\cos b + \sin a \sin b
= \frac12\cdot\frac{\sqrt2}{2} + \frac{\sqrt3}{2}\cdot\frac{\sqrt2}{2}
= \frac{\sqrt2 + \sqrt6}{4},
\]
\[
\sin\frac{\pi}{12} = \sin a\cos b - \cos a \sin b
= \frac{\sqrt3}{2}\cdot\frac{\sqrt2}{2} - \frac12\cdot\frac{\sqrt2}{2}
= \frac{\sqrt6 - \sqrt2}{4}.
\]
\end{solution}

\begin{solution}{exo:g12:trigo:5}
Ambil $s = \sin x$: $2s^2 - s - 1 = (2s + 1)(s - 1) \geq 0$ jika dan hanya jika
$s \leq -\frac12$ atau $s = 1$.

$\sin x = 1$ memberi $x = \frac\pi2$. Untuk $\sin x \leq -\frac12$: pada
lingkaran satuan, itulah busur di bawah garis mendatar $y = -\frac12$, yang di
$\intcc{0}{2\pi}$ menjadi $x \in \intcc{\frac{7\pi}{6}}{\frac{11\pi}{6}}$.
Jadi
\[
S = \left\{\frac{\pi}{2}\right\} \cup
\intcc{\frac{7\pi}{6}}{\frac{11\pi}{6}} .
\]
\end{solution}

\begin{solution}{exo:g12:trigo:6}
$\dfrac{\sin 3x}{x} = 3\,\dfrac{\sin 3x}{3x} \to 3 \cdot 1 = 3$
(komposisi dengan $u = 3x \to 0$).

Dengan memakai $1 - \cos x = \dfrac{\sin^2 x}{1 + \cos x}$:
\[
\frac{1 - \cos x}{x^2}
= \left(\frac{\sin x}{x}\right)^{\!2} \frac{1}{1 + \cos x}
\xrightarrow[x\to0]{} 1 \cdot \frac12 = \frac12 .
\]

Dengan mensubstitusikan $u = \frac1x \to 0^+$:
$x \sin\frac1x = \frac{\sin u}{u} \to 1$.
\end{solution}

\begin{solution}{exo:g12:trigo:7}
$f'(x) = 1 - \sin x \geq 0$ untuk semua $x$, dengan kesamaan tepat ketika
$\sin x = 1$, \emph{yaitu} $x = \frac\pi2 + 2k\pi$. Pada $\intcc{0}{2\pi}$, $f$
naik dari $f(0) = 1$ ke $f(2\pi) = 2\pi + 1$ (nol $f'$ di
$\frac{\pi}{2}$ terpencil, sehingga naiknya bahkan tegas); \omterm{def:g10:functions:extrema}{minimumnya} $1$ di
$0$, \omterm{def:g10:functions:extrema}{maksimumnya} $2\pi + 1$ di $2\pi$. Pada $\R$, $f' \geq 0$ di mana-mana dengan
nol yang terpencil saja, jadi $f$ naik (secara tegas) menurut
\cref{thm:g12:deriv:variations}, walaupun $f'$ bernilai nol di setiap
$\frac\pi2 + 2k\pi$.
\end{solution}

\begin{solution}{exo:g12:trigo:8}
\emph{1.} Misalkan $R = \sqrt{A^2 + B^2} > 0$. Titik
$\left(\frac AR, \frac BR\right)$ terletak pada lingkaran satuan, jadi ada
$\varphi$ dengan $\cos\varphi = \frac AR$ dan $\sin\varphi = \frac BR$. Maka
\[
R\cos(x - \varphi) = R\bigl(\cos x\cos\varphi + \sin x \sin\varphi\bigr)
= A\cos x + B\sin x = f(x).
\]

\emph{2.} Karena daerah hasil $\cos$ adalah $\intcc{-1}{1}$, \omterm{def:g10:functions:extrema}{maksimum} $f$ adalah $R$
dan \omterm{def:g10:functions:extrema}{minimumnya} $-R$. Untuk $\cos x + \sin x$: $A = B = 1$, $R = \sqrt2$,
$\varphi = \frac\pi4$, jadi \omterm{def:g10:algebra:equation}{persamaannya} menjadi
$\sqrt2\cos\left(x - \frac\pi4\right) = 1$, \emph{yaitu}
$\cos\left(x - \frac\pi4\right) = \frac{\sqrt2}{2}$, yang memberi
$x - \frac\pi4 = \pm\frac\pi4 + 2k\pi$: di $\intoc{-\pi}{\pi}$,
$x \in \left\{0, \frac\pi2\right\}$.
\end{solution}

\begin{solution}{exo:g12:trigo:9}
$(\sin)'' = (\cos)' = -\sin$ dan $(\cos)'' = (-\sin)' = -\cos$: keduanya
memenuhi $y'' = -y$.

Misalkan $g = y - a\cos - b\sin$. Maka $g'' = -g$ (kelinearan penurunan),
$g(0) = y(0) - a = 0$ dan $g'(0) = y'(0) - b = 0$. Ambil
$E = g^2 + (g')^2$. Maka
\[
E' = 2g g' + 2g' g'' = 2g g' - 2g' g = 0,
\]
sehingga $E$ tetap dan sama dengan $E(0) = 0$. Jumlah dua kuadrat yang
bernilai nol secara identik memaksa $g = 0$, jadi $y = a\cos + b\sin$.
\end{solution}

\begin{solution}{pb:g12:trigo:1}
\textbf{1.} $0.99833\ldots$ dan $0.99998\ldots$: merayap menuju
$1$.

\textbf{2.} Segitiga $OAM$: alasnya $1$, tingginya $\sin x$, luasnya
$\frac{\sin x}{2}$. Juring $OAM$: bagiannya
$\frac{x}{2\pi}$ dari cakram satuan, luasnya $\frac x2$. Segitiga
$OAT$: alasnya $1$, tingginya $\tan x$, luasnya $\frac{\tan x}{2}$.

\textbf{3.} Segitiga kecil itu duduk di dalam juringnya, dan juringnya di
dalam segitiga besar: $\frac{\sin x}{2} < \frac x2 < \frac{\tan x}{2}$,
yaitu $\sin x < x < \tan x$. Dari $\sin x < x$ diperoleh
$\frac{\sin x}{x} < 1$; dan dari
$x < \tan x = \frac{\sin x}{\cos x}$, dengan mengalikan dengan
$\frac{\cos x}{x} > 0$: $\cos x < \frac{\sin x}{x}$. Ketika
$x \to 0^+$, $\cos x \to 1$, sehingga terapit:
$\frac{\sin x}{x} \to 1$; dan $\frac{\sin x}{x}$ genap, jadi
limit dua sisinya adalah $1$.

\textbf{4.} $\frac{1 - \cos x}{x^2} =
\frac{1 - \cos^2 x}{x^2 (1 + \cos x)} =
\left(\frac{\sin x}{x}\right)^2 \frac{1}{1 + \cos x}
\to 1 \times \frac12$.

\textbf{5.} $\sin'(0) = \lim_{h \to 0} \frac{\sin h - 0}{h} =
1$: seluruh penurunan \omterm{def:g12:trigo:cossin}{sinus} dan \omterm{def:g12:trigo:cossin}{kosinus} mengalir dari
satu limit ini. Diukur dalam derajat, limitnya menjadi
$\frac{\pi}{180}$, dan setiap \omterm{def:g12:deriv:derivative}{turunannya} akan menyeret
tetapan itu: \omterm{def:g11:trigo:radian}{radian} tepat merupakan satuan yang membuat kalkulus
osilasi menjadi bersih.

\textbf{6.} $y' = -R\omega\sin(\omega t - \varphi)$,
$y'' = -R\omega^2\cos(\omega t - \varphi) = -\omega^2 y$.

\textbf{7.} $\omega = 2$. $y(0) = A = 3$;
$y'(0) = 2B = 8$, jadi $B = 4$ dan $y = 3\cos 2t + 4\sin 2t$.

\textbf{8.} $R = \sqrt{9 + 16} = 5$; kala ayunnya
$\frac{2\pi}{2} = \pi$; fasenya
$\varphi = \arctan\frac43 \approx 0.927$, sehingga
$y = 5\cos(2t - 0.927)$.

\textbf{9.} $5\cos(2t - \varphi) = 0$ untuk pertama kalinya ketika
$2t - \varphi = \frac\pi2$, jadi
$t = \frac{\pi/2 + 0.927}{2} \approx 1.25$ sekon.

\textbf{10.} Dengan $\sin\theta \approx \theta$:
$\theta'' = -\frac gL \theta$, yaitu gerak harmonis dengan
$\omega = \sqrt{\frac gL}$ dan kala
$T = \frac{2\pi}{\omega} = 2\pi\sqrt{\frac Lg}$. Untuk
$L = 1$: $T \approx 2.006$ sekon. Kala ayunnya tidak bergantung pada
amplitudonya (untuk ayunan kecil) --- keisokronan Galileo,
kenyataan yang memungkinkan adanya jam bandul.

\textbf{11.} $L = g\left(\frac{T}{2\pi}\right)^2 =
\frac{9.81}{\pi^2} \approx 0.994$ m: bandul sekon itu
kira-kira $6$ mm lebih pendek daripada meter meridian. Akademinya memilih
meridian (gravitasi berubah dari tempat ke tempat, mengkhianati
bandulnya); seandainya $g$ setipis rambut lebih besar, meter kita akan
berdetak.

\textbf{12.} $\cos\left(x + \frac\pi2\right) =
\cos x \cos\frac\pi2 - \sin x \sin\frac\pi2 = -\sin x$;
$\sin\left(x + \frac\pi2\right) = \cos x$. Jadi menurunkan
\omterm{def:g12:trigo:cossin}{sinus} memberi \omterm{def:g12:trigo:cossin}{sinus} yang bergeser seperempat putaran ($\cos$), lalu sekali lagi
($-\sin$), dan sekali lagi, dan pulang dalam empat langkah: penurunan
memutar lingkarannya.

\textbf{13.} Dengan menjumlahkan kedua jabaran $\cos(a \mp b)$:
$\cos(a-b) + \cos(a+b) = 2\cos a\cos b$. Dengan $a = b = x$:
$1 + \cos 2x = 2\cos^2 x$, yaitu
$\cos^2 x = \frac{1 + \cos 2x}{2}$.

\textbf{14.} Dengan $p = \frac{p+q}{2} + \frac{p-q}{2}$ dan
$q$ bayangan cerminnya, lalu \omterm{def:g10:algebra:expand}{menjabarkan} dan menjumlahkan:
$\sin p + \sin q = 2\sin\frac{p+q}{2}\cos\frac{p-q}{2}$. Untuk
$440$ dan $444$ Hz: sebuah nada pada $442$ Hz yang kenyaringannya
dimodulasi oleh $\cos(2\pi \cdot 2t)$ --- kenyaringannya memuncak $4$
kali tiap sekonnya (\omterm{def:g10:numbers:abs}{nilai mutlak} selubungnya): empat layangan
tiap sekonnya, yang lenyap ketika kedua garpu itu sepakat.

\textbf{15.} $\cos 3x = \cos 2x\cos x - \sin 2x \sin x =
(2\cos^2 x - 1)\cos x - 2\sin^2 x\cos x = 4\cos^3 x -
3\cos x$. Di $x = 20^\circ$: $\cos 60^\circ = \frac12$, jadi
$c = \cos 20^\circ$ memenuhi $8c^3 - 6c - 1 = 0$ --- \omterm{def:g10:algebra:equation}{persamaan} pangkat tiga
tanpa penyelesaian berbentuk akar kuadrat: melukis $20^\circ$, yaitu
membagi tiga $60^\circ$, berada di luar jangkauan penggaris dan jangka.

\textbf{16.} $R = \sqrt{9 + 27} = 6$;
$\tan\varphi = \frac{3\sqrt3}{3} = \sqrt3$, jadi
$\varphi = \frac\pi3$ dan $I(t) = 6\cos\left(100\pi t -
\frac\pi3\right)$: puncaknya $6$ ampere, dengan ketinggalan fase $60^\circ$.

\textbf{17.} $\sin^2 t = \frac{1 - \cos 2t}{2}$, jadi kalanya
$\pi$. Kala $\sin 2t$ adalah $\pi$, sedangkan kala $\sin 3t$ adalah
$\frac{2\pi}{3}$: jumlahnya berulang setelah kelipatan persekutuan
terkecilnya, yaitu $2\pi$.

\textbf{18.} Osilasi berkala $1$ di dalam selubung yang menyusut,
$\pm\eu^{-t/10}$: tiap ayunan kira-kira $10\,\%$ lebih rendah. Di
$t = 10$ selubungnya bernilai $\eu^{-1} \approx 0.368$ --- lagi-lagi
tetapan \cref{pb:g12:exp:1}. \omterm{def:g10:algebra:equation}{Persamaan} dengan penyelesaian
teredam seperti itu ($y'' + a y' + b y = 0$) diselesaikan di bab
\omterm{def:g10:algebra:equation}{persamaan} diferensial.

\textbf{19.} Gaya yang menarik kembali sebanding dengan
simpangannya memberi $y'' = -\omega^2 y$; dan
\cref{exo:g12:trigo:9} menunjukkan bahwa penyelesaian \omterm{def:g10:algebra:equation}{persamaan} itu
\emph{tepat} berupa gabungan $\cos$ dan $\sin$ ---
keberadaannya lewat contoh, ketunggalannya lewat latihan itu: getaran
tidak punya pilihan selain menjadi sinusoidal.

\textbf{20.} Satu limit ($\frac{\sin x}{x} \to 1$, lewat
luas yang terapit), dua \omterm{def:g12:deriv:derivative}{turunan} ($\sin' = \cos$,
$\cos' = -\sin$), satu \omterm{def:g10:algebra:equation}{persamaan} ($y'' = -\omega^2 y$), dan sebuah dunia
berisi jam, dawai, arus dan pasang surut. Tanda bintangnya: untuk ayunan
lebar, $T = 4\sqrt{\frac Lg} \cdot
\frac{\pi/2}{M(1, \cos(\theta_0/2))}$ --- yaitu
rata-rata aritmetika--geometri \cref{pb:g12:seq:1} yang menghitung
integral eliptis itu dalam beberapa lelaran saja: catatan buku harian Gauss
dan bandul Galileo, hanya berjarak satu rumus.

\end{solution}
